\chapter{提示}

卡住超过 20 分钟再来看这里。提示按题号排列，从思路到接近答案逐步具体。

\RaggedRight

\section{Part 1}

\begin{description}[style=nextline, leftmargin=1.5em]
  \item[1.1] 紧凑构造函数写在 record 的大括号里面，形式是 \code{public Station { ... }}，没有参数列表。里面可以直接用 \code{name}、\code{availableBikes} 这些参数名。
  \item[1.3] 流水线的顺序：
\begin{plaincode}
stream()
  -> map(...)            转成 NearbyStation
  -> sorted(Comparator.comparingLong(NearbyStation::distanceMeters))
  -> limit(limit)
  -> toList()
\end{plaincode}
\code{limit} 参数检查要放在 Stream 前面，因为 \code{Stream.limit} 收到负数时抛的异常信息不清楚。
  \item[1.4] 可以连着写几个 \code{filter}。比如“如果打开了 requireBikes，就只保留有车的”，可以写成 \code{.filter(s -> !requireBikes || s.availableBikes() > 0)}。想一想为什么这样写是对的。
\end{description}

\section{Part 2}

\begin{description}[style=nextline, leftmargin=1.5em]
  \item[2.1 报错 No qualifying bean of type AppProperties] 说明容器里没有 \code{AppProperties} 这个 Bean。检查两处：\code{AppProperties} 上有没有 \code{@ConfigurationProperties(prefix = "app")}；\code{AppConfig} 上有没有 \code{@Configuration} 和 \code{@EnableConfigurationProperties(AppProperties.class)}。
  \item[2.2 嵌套对象是 null] 在 \code{AppProperties} 的 \code{jcdecaux} 和 \code{scraper} 两个参数前面加 \code{@DefaultValue}。
\end{description}

\section{Part 3}

\begin{description}[style=nextline, leftmargin=1.5em]
  \item[3.1 某个字段是 0 或 null] 说明 Jackson 没找到对应的 JSON 字段。检查 \code{@JsonProperty} 里的名字是否和 JSON 完全一致，以及 import 的是不是 \code{com.fasterxml.jackson.annotation.JsonProperty}。
  \item[3.3 报错 No further requests expected 或 URL 不匹配] 把报错信息里实际请求的 URL 和期望的 URL 逐个字符对比。常见问题：\code{path} 少了开头的 \code{/}，或者两个查询参数的顺序反了。
  \item[3.3 构造函数里 restClient 赋值] \code{this.restClient = builder.baseUrl(props.baseUrl()).build();}
\end{description}

\section{Part 4}

\begin{description}[style=nextline, leftmargin=1.5em]
  \item[4.4 派生查询的方法名] 拆开来看：\code{findFirst} + \code{By} + \code{IdNumber}（条件：id.number 等于参数）+ \code{OrderBy} + \code{IdLastUpdate} + \code{Desc}。
  \item[4.5 JPQL 写不出来] 照着知识点里“每个作者最新的书”改：实体是 \code{AvailabilityEntity}，分组字段是 \code{a.id.number}，比较的字段是 \code{a.id.lastUpdate}。
\end{description}

\section{Part 5}

\begin{description}[style=nextline, leftmargin=1.5em]
  \item[5.1 不知道从哪里开始] 一个 for 循环，两个计数器 \code{stationsSaved} 和 \code{availabilityInserted}。站点部分照着知识点里的 upsert 写；可用量部分：先用 \code{dto.lastUpdateInstant()} 判断是否为 null（是 null 就 \code{continue}），再 \code{new AvailabilityId(dto.number(), lastUpdate)}，用 \code{existsById} 判断。
  \item[5.3 @Scheduled 的写法] \code{@Scheduled(fixedDelayString = "${app.scraper.interval:5m}")}
\end{description}

\section{Part 6}

\begin{description}[style=nextline, leftmargin=1.5em]
  \item[6.1a] 第一步，把 \code{findLatestForAllStations()} 的结果转成 Map：
\begin{java}
Collectors.toMap(a -> a.getId().number(), Function.identity())
\end{java}
第二步：\code{stationRepository.findAll(Sort.by("number"))} 转成 Stream，\code{map} 里调用 \code{toStation(s, map.get(s.getNumber()))}。
  \item[6.1b] 先找站点，找不到就抛异常；再找最新的可用量，找不到就用 null：
\begin{java}
stationRepository.findById(number)
        .orElseThrow(() -> new StationNotFoundException(number));
availabilityRepository.findLatest(number).orElse(null);
\end{java}
  \item[6.3 404 测试过了，400 测试的 Content-Type 不对] 检查 \code{GlobalExceptionHandler} 是否继承了 \code{ResponseEntityExceptionHandler}。
  \item[FinalApplicationTest 失败] 找报错里 \code{Caused by:} 开头的最后一行，通常是 \code{No qualifying bean of type 'XXX'}，说明 XXX 这个类没有被 Spring 管理：检查它上面有没有 \code{@Service} 或 \code{@Component}。
\end{description}
